\documentclass[11pt]{article}
\usepackage[T1]{fontenc}
\usepackage{lmodern}
\usepackage{amsmath,amssymb}
\usepackage{booktabs}
\usepackage{array}
\usepackage[a4paper,margin=2.4cm]{geometry}
\usepackage[hidelinks]{hyperref}
\setlength{\parskip}{5pt}
\title{\vspace{-1.6cm}\Large Basics and Definitions\\[2pt]
\large Everything you need before the research plan, written for a reader who has finished class ten}
\author{Dr.\ Vipin \\ Gurugram University}
\date{30 October 2026}
\begin{document}
\maketitle
\vspace{-6mm}

\section{How to use this booklet}

This booklet has one job: to make the research plan readable without a dictionary. The plan talks about
guided waves in a plate, about strain and stress, about electric fields in a material that conducts a
little, about heat that remembers, and about a size effect that appears only when the plate is thin. Each of
those ideas can be explained in ordinary language, and that is what the following pages do. Where an idea
needs algebra, the algebra is written out slowly and no step is skipped.

You do not have to read the booklet in one sitting. Sections two to nine are about waves and are the
foundation. Sections ten and eleven are about the complex wavenumber, which sounds frightening and is not.
Sections twelve to fourteen are about how solids carry force. Sections fifteen to nineteen are about the
four couplings in our model. Section twenty is about what a model is allowed to assume. Section twenty one
is a table of every symbol in the plan, section twenty two is a glossary, and section twenty three tells you
how to test your own work when nobody is looking over your shoulder.

Nothing here contains a boundary condition, a solution method, or a number to compare with. Those are your
job, and they are the reason the plan exists.

\section{What a wave is}

A wave is a disturbance that travels. The disturbance moves; the material need not. If you drop a stone into
a pond, the rings spread outwards, but the water itself mostly bobs up and down in place. If you speak, the
air carries the disturbance to your friend, but the air does not travel from your mouth to his ear; it
vibrates and passes the message on. A wave, then, is a pattern of motion that hands itself from one piece of
material to the next.

Two things are worth fixing in the mind at the start. The first is that a wave carries energy without
carrying material. That is why a sound can break a glass without any of the air in the room ending up inside
the glass. The second is that the speed of a wave is a property of the material, not of how hard you hit it.
Hit a table harder and the sound is louder, not faster. This is true for ordinary waves and it is true for
the waves in our plate, with one important exception that we come to in section six.

A wave has an amplitude, which is how big the disturbance is, and a shape, which is how the disturbance is
distributed in space. For the waves in this project the shape is a smooth curve, either a sine or a
combination of sines, and that is not a simplification for convenience. Any smooth repeating shape can be
built out of sines, so if we understand what happens to a single sine we understand what happens to
everything.

The disturbance in a solid can be of two kinds. It can be a push and pull along the direction the wave
travels, which is called a longitudinal wave, and it can be a sideways sliding across the direction of
travel, which is called a transverse wave. Sound in air is longitudinal. Light is transverse. Earthquakes
carry both, which is why a seismograph shows two arrivals for one event. A solid can carry both, and that is
the root of everything that follows: in a plate, the two kinds get mixed together by the free surfaces.

\section{Frequency and period}

Frequency is how many times per second the disturbance repeats. Its unit is the hertz, written Hz, and one
hertz is one cycle per second. A tuning fork at 440 Hz pushes the air 440 times every second, which is why
the note sounds the way it does.

The period is the time for one cycle, and it is the reciprocal of the frequency. If the frequency is 440 Hz
the period is one four hundred and fortieth of a second, a little over two milliseconds. This reciprocal
relationship is worth remembering because different problems are easier in different units. A fast
oscillation has a high frequency and a short period; a slow one has a low frequency and a long period.

There is a second measure of frequency, the angular frequency, written with the Greek letter omega. It is
the frequency multiplied by two pi. The reason for the multiplication is that a full cycle is a full turn,
and a full turn is two pi radians. If you measure angles in radians rather than in degrees, the arithmetic of
oscillations becomes clean, and the clean arithmetic is worth the extra symbol. Whenever this project says
omega, read it as frequency in radians per second.

Frequency also has a physical meaning that goes beyond bookkeeping. It decides how fast the material is
being asked to change, and every material has rates at which it can respond comfortably and rates at which it
cannot. A material that follows a slow push faithfully may lag behind a fast one. When it lags, energy is
left behind as heat, and the wave weakens. This is the origin of most of the attenuation in the paper, and
we will come back to it.

\section{Wavelength}

Wavelength is the distance between two neighbouring crests, or more generally between two points that are
doing the same thing at the same time. Its unit is the metre or a fraction of it.

Wavelength matters because a wave only notices objects that are comparable with its own wavelength. Sound
with a wavelength of a metre passes around a doorway of half a metre without much trouble. Sound with a
wavelength of a centimetre is blocked and reflected, which is why a high note from one room is harder to hear
in the next than a low one. In a plate, the wavelength relative to the thickness decides which kind of motion
is possible, and this is the central geometric fact of the whole project.

There is a simple relation between the three quantities so far. The wave moves one wavelength in one period.
So the speed is the wavelength divided by the period, and since the period is the reciprocal of the
frequency, the speed is the wavelength multiplied by the frequency. If a wave has a wavelength of one metre
and a frequency of two hertz, it moves two metres per second. Nothing more complicated than that is needed
to follow the paper.

\section{The wavenumber}

The wavenumber is a way of measuring how quickly the wave pattern repeats in space, rather than how quickly
it repeats in time. Frequency does the counting in time; the wavenumber does the counting in space. Its
symbol is k and its unit is the reciprocal of a metre.

The most useful definition is not the number of crests per metre but the phase per metre. One full cycle is
two pi radians of phase, so a wave whose wavelength is lambda has a wavenumber of two pi divided by lambda.
If you keep the two pi in the definition, then a wave travelling in the positive direction is written as a
cosine with the argument k x minus omega t, and every formula that follows comes out without stray factors
of two pi. That is the whole reason for the definition, and it is a good enough reason.

Here is a way to feel the difference between the two counts. Stand on a beach and count waves arriving per
minute; that is frequency. Walk along the beach and count how many wave crests fit into a hundred metres;
that is wavelength, and its reciprocal is the wavenumber up to the factor of two pi. Both describe the same
train of waves, one in time and one in space.

For our plate the wavenumber is the important one, because the plate is infinite along the direction of
travel and finite across it. The wave is free to choose a wavenumber along the plate, and the plate then
decides which through-thickness shapes are allowed at that wavenumber. The paper is, in the end, a
calculation of that decision.

\section{Phase velocity, and why it is not the only speed}

Phase velocity is the speed at which a crest moves. It is the frequency divided by the wavenumber, or
equivalently the wavelength multiplied by the frequency, which is the same statement as before.

If phase velocity were the same for every frequency, waves would be simple. They are not. In most real
media the phase velocity depends on frequency, and the dependence is called dispersion. A glass prism
separates white light into colours precisely because the phase velocity of light in glass depends on
frequency. Sound in a plate behaves in the same way: the speed of a guided wave depends on how fast it is
wobbling.

There is a second speed, the group velocity, which is the speed at which a package of waves travels and the
speed at which energy moves. For a single pure sine the two are the same, but a real pulse is built out of
many frequencies, and if each frequency travels at its own speed the pulse changes shape as it goes. That is
why a hand clap heard across a valley arrives as a smear rather than a click. The paper reports phase
velocity, because that is the quantity a dispersion curve is drawn with and the quantity a device is designed
around, and it mentions group velocity where the energy argument needs it.

Dispersion is a nuisance for a communication engineer and a gift for us. A measured dispersion curve is a
fingerprint. If you know how phase velocity varies with frequency for a given plate, you can work backwards
and say something about the material, the thickness, or the presence of damage. This is the basis of a large
industry of non destructive testing, and it is the reason anybody cares about our model.

\section{A plate, and why waves get trapped in it}

A plate is a solid whose thickness is much smaller than its length and width. A sheet of window glass is a
plate. A coin is a short cylinder, which is not the same thing. In this project the plate has a thickness
written as two h, so the two faces sit at plus h and minus h, and it is infinite in the two other
directions. Infinite is a mathematician's word for long enough that the ends do not matter.

Now send a wave into the plate at an angle. It travels through the material until it meets a face. A face is
free, which means nothing is pressing on it, so the wave cannot push its way out into the air in a useful
way; almost all of it reflects back into the plate. It then travels across the thickness, meets the other
face, reflects again, and so on. The wave is trapped, and while it is trapped it creeps along the plate.

This is not a small effect, it is the whole effect. A wave that is bouncing between two faces is not
travelling straight; it is travelling in a zigzag, and the zigzag has a forward component. That forward
component is the guide wavelength and the speed along the plate is the phase velocity you measure with a
transducer at the surface. The steeper the zigzag, the slower the apparent speed along the plate, because the
wave is spending most of its effort going up and down rather than forward.

Two consequences follow at once, and both matter for the paper. First, the thickness enters the problem
directly, because the zigzag path length is set by it. Second, the two faces talk to each other: the wave
returns to a face with a phase that depends on the thickness, the wavelength, and the angle. When the return
phase is just right, the two reflections reinforce, and a mode exists. When it is not, the mode cannot exist,
and the amplitude dies away. That is where the discrete set of modes comes from.

\section{Modes, and the two families}

Because the plate can only support the zigzags whose return phase fits, the waves in a plate come in a
discrete list rather than a continuum. Each entry in the list is called a mode. A mode is a definite
through-thickness shape together with a dispersion relation of its own. You can think of it as a standing
pattern across the thickness that is carried along the plate.

The modes fall into two families, and the reason is symmetry. The plate looks the same if you turn it upside
down, so its solutions must come in pairs related by that turn.

In the symmetric family the two faces move towards each other and away from each other in step, so the middle
of the plate moves along the plate and does not move across it. The lowest member, called S zero, is the one
that stretches and squeezes the plate like a concertina. In the antisymmetric family the two faces move in
the same direction at the same time, so the plate bends. The lowest member, called A zero, is the bending
wave of a plate and it is the one a plate feels when you tap it.

At low frequency only these two exist. As the frequency rises, more modes appear, in the order S one, A one,
and so on, and each new mode appears at a frequency called its cutoff. The two families are independent only
when the material is simple. In our materials the electric field and the temperature couple the families
together at the faces, so in general a wave that starts as symmetric picks up a small antisymmetric part as
it travels. Where that mixing is large enough to matter is one of the questions the paper has to answer.

\section{Cutoff, dispersion curves, and the product of frequency and thickness}

The cutoff of a mode is the lowest frequency at which it can exist. Below cutoff the mode is not a travelling
wave at all; it dies away exponentially from wherever it was created, and it is called evanescent. The
physical reason is the zigzag: at low frequency the wavelength is long compared with the thickness, so a wave
cannot fit a half wavelength across the plate and still come back in step. Only the two fundamental modes
avoid this restriction, and that is why a thin plate at low frequency carries only S zero and A zero.

A dispersion curve is a graph of phase velocity against frequency, one curve per mode. On such a graph the
cutoffs appear as points where a curve begins. The curves of the higher modes rise steeply after their cutoff
and then flatten out at high frequency towards the speed of a wave on the surface of a half space, called the
Rayleigh speed. That flattening is a useful test: any calculation that does not tend to the Rayleigh speed at
high frequency has something wrong with it, and section twenty three says so again.

One practical detail saves a great deal of confusion. A plate of twice the thickness at half the frequency
behaves in almost the same way as the original, because the wave does not know about frequency and thickness
separately; it knows about their product. So dispersion curves are drawn with frequency multiplied by
thickness along the horizontal axis, with unit hertz times metre. On such an axis all plates of the same
material fall on the same set of curves, and thickness stops being a nuisance.

\section{Complex numbers, from zero}

The phrase complex number frightens people for no good reason. A complex number is an ordinary number with
one extra idea attached, and the extra idea is this: there is a number whose square is minus one.

Ordinary numbers do not have that property. Anything real squared is positive, or zero if it is zero. So the
equation x squared equals minus one has no solution among the numbers you have used so far. Mathematicians
solved the problem by inventing a solution and calling it i, with the rule that i squared is minus one. It
looks like a trick, and for a long time it was treated as one. It is not a trick: it is a bookkeeping device
for two quantities at once, and once you have it, a great many things that would otherwise need two equations
need only one.

A complex number is written as a plus b i, where a and b are ordinary numbers. The first part is called the
real part and the second is called the imaginary part. The word imaginary is a historical accident; there is
nothing imaginary about the second number. If a is the length of a step to the east and b is the length of a
step to the north, then a plus b i is simply the position after both steps. The complex number is a pair of
numbers written compactly.

Complex numbers add easily: you add the real parts and add the imaginary parts separately, exactly as you add
east and north steps separately. They multiply by the usual rule with i squared replaced by minus one. For
example, (one plus i) times (one plus i) is one plus two i plus i squared, which is two i. The real part
cancelled.

The most useful fact about complex numbers is the size. The size, or modulus, of a plus b i is the square root
of a squared plus b squared. That is Pythagoras, and it is the length of the step to the corner. The size is
written with two vertical bars, and it is always positive, whatever the signs of a and b.

The second most useful fact is that a complex number can be written in polar form, as a size times a phase.
Picture the step to the corner. Its length is the size and its angle from the east direction is the phase.
The phase is an angle in radians. When you multiply two complex numbers, their sizes multiply and their
phases add. That single sentence carries a great deal of physics, because a wave is exactly a size and a
phase, and a wave travelling through a material is exactly a size that shrinks and a phase that turns.

Now the exponential. There is a famous relation, due to Euler, which says that the exponential of i times an
angle equals a cosine plus i times a sine. A cosine and a sine are the real and imaginary parts of the same
turning arrow. So one can write a wave as a size times the exponential of i times something, and the size is
its amplitude while the something is its phase. This notation is used everywhere in acoustics because it
turns the trigonometry of waves into the arithmetic of exponents, where multiplication is easy and
differentiation gives back the same function.

Two operations look after themselves in this language. Differentiating with respect to time multiplies by
minus i omega. Differentiating twice multiplies by minus omega squared. Differentiating with respect to
position multiplies by i k. That is why the equations in the plan look so short: a symbol like omega squared
in front of a term is a second time derivative in disguise, and a symbol like k squared is a second space
derivative in disguise. If you can read that, you can read most of the paper.

\section{The complex wavenumber, and what attenuation means}

A wave moving in the positive direction along a plate can be written as a size times the exponential of i
times the quantity k x minus omega t. If k is an ordinary positive number, this describes a wave of constant
amplitude travelling forever. That is unphysical, because real waves weaken.

So we let k be complex. Write k as k real plus i k imaginary. Put that into the exponential and split it:

the exponential of i k x equals the exponential of i k real x times the exponential of minus k imaginary x.

The first factor turns the phase and the second shrinks the size. So the real part of k decides the
wavelength, through the usual relation that the wavelength is two pi divided by the real part, and the
imaginary part decides how quickly the wave dies away with distance.

The convention in this project is that a positive imaginary part means decay in the positive direction. That
is a choice, like choosing to measure height upwards; the physics is the same either way, but a paper must
say which way it chose, because a reader who assumes the other convention will read a growth where the
author wrote a decay. Every figure in this project therefore states the sign convention once, in the caption.

Attenuation is the rate at which the amplitude falls. If the amplitude falls as the exponential of minus k
imaginary x, then a natural measure of loss is k imaginary itself, in reciprocal metres, which is called
nepers per metre. Nepers are like radians: they are a ratio disguised as a unit. Engineers often prefer
decibels, and the conversion factor is a little over eight and a half decibels per neper, because a decibel
is defined on a logarithmic scale in base ten while a neper is defined in base e.

The ratio of the imaginary part to the real part is the quantity that tells you whether a mode is usable. It
is the fraction of amplitude lost per radian of phase, which is the same as the fraction lost per unit of
distance measured in wavelengths. A mode with a small ratio travels many wavelengths; a mode with a ratio
near one disappears within a wavelength and is useless for a delay line. When the plan says that a certain
effect is too small to see, it usually means that the change in this ratio is smaller than the accuracy with
which it could be measured.

Finally, a word about where attenuation comes from in our problem. Some of it is leakage, some is
thermoelastic loss, and some is the drag of the moving charge carriers on the crystal. Leakage depends on the
surroundings, and we exclude it by keeping the faces free. Thermoelastic loss is a property of every solid
and is small. Carrier drag is the new ingredient in this paper, it can be switched on and off by a bias
voltage, and it can even be made negative, which means amplification. That last possibility is the reason the
subject is interesting for devices.

\section{Stress and strain}

Hold one end of a rubber band and pull the other. The band gets longer. The change in length divided by the
original length is the strain. Strain is dimensionless, because it is a length divided by a length, and it
is usually quoted as a small number or as a percentage. A strain of one per cent means the material got one
per cent longer.

While you are pulling, the band pulls back. The force it transmits, divided by the cross section over which
it acts, is the stress. Stress has the unit of force per area, which is the pascal, and one pascal is one
newton on one square metre. Atmospheric pressure is about a hundred thousand pascals, and steel can carry
hundreds of millions of pascals before it yields, which is why the pascal is a small unit and why engineers
prefer to work in megapascals.

Two kinds of stress matter in a plate. A normal stress pulls two faces apart or pushes them together, and the
rubber band carries a normal stress. A shear stress slides one layer over the next, and the same rubber band
carries a shear stress if you twist it instead of pulling it. A solid can carry both at the same time, and in
a plate it usually does. The reason is again the zigzag: a wave that bounces between the faces converts
stretching into sliding and back, so at any point the material is doing both.

Stress and strain are not single numbers, they are collections. At a point in a solid there are three
directions, and for each direction of facing there is a stress with three components. That is nine numbers,
and for a symmetric object only six are independent. When the plan writes an index like i j, it is naming one
of those components. You do not need to hold all nine in your head. You need to know that the notation is
counting them, and that a material property attached to two indices is a number that relates a stress to a
strain in a particular pair of directions.

There is one more idea, and it is the one that makes the size effect possible. Strain can be non uniform. The
top of a bent beam is stretched more than the middle. The rate at which the strain changes from point to point
is the strain gradient, and it has the strange unit of one over length. In an ordinary theory of elasticity
the energy depends only on the strain; the strain gradient costs nothing. In the theory we use, it costs a
little, and that little is what makes a thin plate stiffer than scaling would predict.

\section{Hooke's law, and how a material answers}

Hooke's law is the statement that the response is proportional to the cause, for small causes. Hang a weight
on a spring and it stretches. Double the weight and it stretches twice as much, provided you do not
overstretch it. The constant of proportionality is the stiffness. This is the whole content of Hooke's law,
and the reason a whole branch of engineering exists is that the same simple idea works for solids with a
little more bookkeeping.

In a solid the bookkeeping is this. A strain in one direction produces stresses in several directions.
Pull a strip of rubber, and it gets longer but also thinner; the thinning is a strain in the two directions
across the pull. The ratio of the sideways contraction to the lengthwise extension is a material property
called Poisson's ratio. A value near one half means the material keeps its volume when stretched, which is
what rubber does. A value near zero means it hardly narrows at all, which is what cork does, and that is why
cork is used in bottle stoppers.

The full version of Hooke's law for an elastic solid is a table of numbers that relates six stresses to six
strains. For an isotropic material, which looks the same in every direction, the table collapses to two
numbers. These are often chosen as the pair called the Lam\'e constants, one for stiffness in general and one
for resistance to shear. It is convenient to define two speeds from them: a speed for push pull waves and a
speed for shear waves. The push pull speed is always the larger of the two, because a material resists being
squeezed more than it resists being sheared. In the model we write the table as a symbol with four indices,
which is intimidating until you notice that in the isotropic case it only contains two different numbers.

Our materials are not isotropic. A piezoelectric crystal has a preferred axis and looks different along it
than across it; the number of independent constants rises, but only to five or so for the symmetry class we
use. A laminated composite is not even a single material, it is a stack, and its effective stiffness is an
average that depends on how the layers are arranged. The paper deals with this by computing effective
constants once and using them throughout, which is the standard engineering compromise and is stated as an
assumption.

\section{Balance laws, and what a divergence means}

Physics has a small number of conservation statements and a great many consequences. Three matter here.

Mass is conserved. What flows into a small box either stays in it or comes out; nothing is created. Momentum
is conserved. The force on a box equals the rate at which its momentum changes, which for a solid means the
weight of the material times its acceleration. Energy is conserved, and for a solid at constant temperature
this reduces to a balance between stored elastic energy, kinetic energy, and the work done by outside forces.

A balance law is usually written with a divergence, and the word alarms people. A divergence is a net
outflow. Imagine a small box, and imagine a flow of something through it. If more flows out through the six
faces than flows in, the amount inside the box falls, and the divergence is positive. If more flows in, the
divergence is negative and the amount inside rises. That is all it means. When the plan writes the divergence
of a stress equal to density times acceleration, it is saying that if the stresses on the faces of a small
box do not balance, the box accelerates. When it writes the divergence of the electric displacement equal to
minus the charge density, it is saying that electric displacement lines start and end on electric charge.

A wave equation comes from combining a balance law with a material law. The material law says how stress
depends on strain; the balance law says how stress accelerates the material. Put the first into the second
and you get an equation in which the acceleration of the displacement is proportional to its curvature in
space. That equation has, as its solutions, waves whose speed is the square root of the ratio of the
stiffness to the density. This is the origin of every wave speed in the paper, and it is worth seeing once,
because after that the paper just does the same thing with more fields and more bookkeeping.

\section{Electric fields, and what a potential is}

An electric charge feels a force if it is near other charges. It is convenient to describe that force in two
steps. First, say that the region around a charge is filled with an electric field. Second, say that another
charge placed in that field feels a force equal to its own charge times the field. The field is the force per
unit charge, and its unit is the volt per metre.

The potential is a second convenience, and it is the one beginners find hardest. The potential is the energy
per unit charge. Moving a charge from one place to another costs energy, and the cost divided by the charge is
the difference in potential, which is what we call voltage. A nine volt battery is a promise that every
coulomb of charge taken through it will release nine joules of energy.

The relation between field and potential is a gradient, with a minus sign. The field points in the direction
in which the potential falls fastest. Water flows downhill; charge drifts towards lower potential. The minus
sign in the definition is there so that the two statements agree. In an equation it looks like the electric
field being minus the derivative of the potential, and in words it is the sentence you just read.

There is one more idea to carry, and it is Gauss's law. Electric field lines begin on positive charge and end
on negative charge. If you draw a closed surface around some charge, the number of field lines poking out
through it, counted with sign, equals the charge inside. In a material, the correct quantity to count is not
the field but the electric displacement, which is the field multiplied by the permittivity of the material.
The reason for introducing the displacement is that the material itself polarises: its atoms stretch into
little dipoles and cancel part of the charge you put in. The displacement counts only the charge you put in,
which is the charge you control.

For our plate this becomes a small bookkeeping question at the faces. On the faces we must say whether the
displacement is continuous with the field outside, or whether it vanishes because the face is covered with a
metal electrode that carries away any charge. Both conditions are used in practice and they give slightly
different dispersion curves. The choice is a modelling decision, and it belongs to you, not to this booklet.

\section{What makes a semiconductor different}

A metal conducts because its electrons are free to move. An insulator does not conduct because its electrons
are held in bonds. A semiconductor sits between the two, and it is the most interesting of the three because
its conductivity can be changed by adding impurities, by shining light on it, or by heating it.

The word for the mobile electrons is carriers. Their number per unit volume is the carrier density. In a pure
semiconductor at room temperature the number is modest; a doped one can have a thousand times more. When the
material is at rest, the carriers are spread evenly and every point is electrically neutral, because each
mobile electron is balanced by a fixed positive ion of the lattice behind it.

Now squeeze the material. The piezoelectric effect produces a voltage across it, which is to say an electric
field. The carriers feel that field and move. They move until their own charge distribution has cancelled
most of the field, and the result is a small region near each face where the charge is not neutral. This
region is called a depletion layer, and it is the working part of a diode.

When a wave passes, the field it produces is not steady. It rises and falls, and the carriers try to follow.
They do not follow perfectly, because they need time. Two processes take that time. Diffusion is the spreading
of carriers from crowded places to empty ones, as a drop of ink spreads in water. Drift is the steady push of
the field on the carriers, and the measure of how readily they respond is the mobility. The product of
mobility and field is a velocity, and that velocity is called the drift velocity.

The point that makes our problem interesting is that the carriers are slow compared with the wave. A wave at
a hundred megahertz in a crystal might travel at three thousand metres per second and take three hundred
nanoseconds for a centimetre. The carriers may need microseconds to redistribute. So the wave hands the
carriers a push, the carriers slowly respond, and the response arrives out of step with the wave that caused
it. An out of step response extracts energy, and that is attenuation. Make the carriers run faster, by
increasing the applied field, and the response can come back into step and even overtake the wave, which is
how an acoustic amplifier works.

There is one further subtlety, which is the reason the plan has two coefficients rather than one. If the
plate is not at a uniform temperature, the carriers feel the temperature gradient as well as the electric
field. They drift down the temperature gradient, from hot to cold, exactly as they drift down a voltage
gradient. This is called thermodiffusion, and it is the same physics as the Seebeck effect in a thermocouple.

\section{Seebeck, Peltier, and reciprocity}

A thermocouple is two wires of different metals joined at both ends. Heat one junction and a voltage appears
at the other. That is the Seebeck effect, and the coefficient that measures it, volts per kelvin, is the
Seebeck coefficient. For a material whose carriers are electrons the coefficient is negative, because the
electrons migrate towards the cold end and leave the hot end positive. For a material whose carriers are
holes the coefficient is positive.

The reverse effect exists. Drive a current through a junction and heat is absorbed at one side and released
at the other. That is the Peltier effect, and its coefficient is in volts, because it is a heat per unit
charge. A Peltier cooler is a small refrigerator with no moving parts, and it is built from many junctions in
series electrically and in parallel thermally.

The two coefficients are not independent. Kelvin showed, from thermodynamics alone, that the Peltier
coefficient equals the temperature times the Seebeck coefficient. Later, with the machinery of irreversible
thermodynamics, Onsager showed why: the two effects are a pair of cross flows, and the coefficients that
couple the flows to the forces must be equal. In the language of the plan, the constant that appears in the
carrier equation and the constant that appears in the heat equation are the same number wearing two hats.

This is a gift for a paper. It means we introduce one new material property, the Seebeck coefficient, and get
both the carrier effect and the heat effect from it. A paper that introduced two independent coefficients
would need two independent measurements to defend them, and nobody would believe either.

There is a sign convention trap here. Two
signs must be kept straight: the sign of the Seebeck coefficient, negative for electrons, and the side of the
equation on which the reciprocal term is written. In our heat equation the source terms appear inside a
bracket with the opposite sign to their contribution, so a positive heating appears as a negative entry. If you ever find yourself unsure, ask which way the carriers are
going and whether the spot they are leaving is getting warmer or cooler, and let the sign follow the physics.

\section{Heat: conduction, diffusion, and a memory}

Heat flows from hot to cold. The flux is proportional to the temperature gradient, with a minus sign, and the
constant is the thermal conductivity. That is Fourier's law, and for steady problems it is all you need.

For time varying problems, Fourier's law is combined with energy conservation, and the result is a diffusion
equation for temperature. The equation says that a point warms if it is colder than its neighbours on average,
which is another way of saying the curvature of the temperature is positive. The rate at which it happens is
set by a single number, the thermal diffusivity, which is the conductivity divided by the density and the
specific heat.

The diffusion equation has an uncomfortable property. Its solutions allow a temperature change at one point
to be felt everywhere at once, however small the effect far away. For most engineering purposes this does not
matter, because the effect is negligible. For a wave oscillating at a very high frequency it does matter,
because the equation predicts that the loss grows without limit as the frequency rises. Nature does not
behave that way, and the divergence is a sign that the theory is being asked a question outside its range.

The repair, due to Maxwell in another context and later to Cattaneo and to Vernotte, is to give the heat flux
a short memory. The flux does not respond instantly to the gradient; it relaxes towards the gradient over a
time called the relaxation time. One such memory gives a hyperbolic heat equation with waves of heat in it.
Two memories, one for the flux and one for the gradient of the flux, give the version used in this project,
which was introduced by Moore, Gibson and Thompson in 1994 and is called the Moore-Gibson-Thompson theory,
or MGT for short.

What does MGT do to a wave in a solid? At low frequency, nothing worth measuring: the relaxation time
multiplied by the frequency is small, and the classical answer returns. At high frequency, the relaxation
takes over, the heat flux lags behind the temperature gradient, and the attenuation stops growing and levels
off at a plateau. The plateau is the physical difference between the classical and the MGT descriptions, and
it is what the paper will look for.

The base form we use has one lag and one time derivative cancelled, which is the simplest member of the
family and the one that has been implemented most often in the recent literature. It has two limits that can
be tested: setting the relaxation time to zero returns Fourier, and letting it become very large gives an
almost adiabatic material. Those two limits bracket the answer, and section twenty three recommends the
bracket as a test.

One more number will keep appearing. If the plate has thickness two h, the time for heat to diffuse across it
is of order the thickness squared divided by the diffusivity. The ratio of the relaxation time to that
diffusion time decides whether the memory is visible for a plate of that thickness. Because thick plates
diffuse slowly and thin plates diffuse quickly, the visibility depends on thickness, and that is why one of
the figures in the plan is a map of thickness against relaxation time rather than a single curve.

\section{The size effect in words}

At the scale of a micrometre, and certainly at the scale of a tenth of a micrometre, a plate does not behave
like a shrunken version of a thick plate. Two separate ideas are used to describe the difference, and both
are needed because experiment shows both kinds of deviation.

The first idea is nonlocal elasticity. The stress at a point is allowed to depend on the strain in a small
neighbourhood around that point as well as on the strain at the point. The image is a row of springs linked by
weak cross springs: a displacement in one place is felt a little further along. Nonlocality softens the
response, so a beam bends more than the classical theory says.

The second idea is strain gradient elasticity. The energy is allowed to depend on how rapidly the strain
changes from point to point, in addition to how large it is. The image is a sheet of rubber on a table: a
region of uniform stretch and a region of rapidly changing stretch can have the same strain at their centres,
but the second costs more energy. The gradient stiffens the response.

The version used in this project combines the two, and it is due to Lim, Zhang and Reddy. One length scale
accompanies each idea, one for the nonlocality and one for the gradient, and the two are in competition: one
makes the plate softer, the other makes it stiffer. Both lengths are tiny, and both are written as a ratio to
the thickness, because a length alone tells you nothing and a ratio to the geometry tells you everything.

There is a price to pay, and it is the sentence that closes the model section of your research plan. When a
theory contains a second derivative of the strain, the equations rise in order. Higher order equations need
more conditions at the boundaries. A classical elastic plate needs three conditions at each face; a plate
described this way needs more, and which extra conditions to impose is not settled in the literature. That is
why the plan hands that decision to you instead of pretending it is part of the model.

\section{What a model is allowed to assume}

A model is not a description of a piece of the world. It is a set of statements that are true enough, in a
range, to answer a question. Knowing the range is part of knowing the model, and a paper that hides the range
is lying quietly.

Our model assumes that the material is a continuum, which is to say that we never look at individual atoms.
It assumes that the response is linear, so that doubling the cause doubles the effect. It assumes that the
properties do not change with temperature, which is a strong assumption given that one of the fields is the
temperature. It assumes that the two length scales of the size effect are material constants and not fitting
parameters. It assumes that the excess carrier density is small compared with the equilibrium density, so
that products of small quantities can be dropped. It assumes that the surfaces are free of any pressure and
free of charge, which is an idealisation of a real surface with an oxide layer on it.

Each assumption is a place where a referee can push, and each should appear in the paper with a sentence of
justification. The good news is that assumptions can be tested. If a result does not change when an
assumption is relaxed, the assumption was not load bearing; if it does, the reader deserves to know which way.

\section{Table of symbols used in the plan}

\begin{center}
\small
\begin{tabular}{p{1.9cm}p{6.0cm}p{2.4cm}p{2.6cm}}
\toprule
symbol & meaning in words & unit & rough size\\
\midrule
$u_x,u_z$ & displacement of a material point along the plate and across it & m & $10^{-10}$ to $10^{-8}$\\
$\sigma_{ij},t_{ij}$ & local stress, total (nonlocal) stress & Pa & $10^{5}$ to $10^{8}$\\
$\varepsilon_{ij}$ & strain, change in length per unit length & none & $10^{-6}$ to $10^{-4}$\\
$\theta$ & temperature rise above the reference value & K & below 1\\
$T_0$ & reference temperature & K & about 300\\
$\Delta n$ & excess carrier density over equilibrium & m$^{-3}$ & $10^{16}$ to $10^{22}$\\
$n_0$ & equilibrium carrier density & m$^{-3}$ & $10^{20}$ to $10^{24}$\\
$q$ & magnitude of the charge of one carrier & C & $1.6\times10^{-19}$\\
$b_n$ & mobility, drift speed per unit field & m$^2$V$^{-1}$s$^{-1}$ & $10^{-3}$ to $10^{-1}$\\
$D_n$ & carrier diffusion coefficient & m$^2$s$^{-1}$ & $10^{-5}$ to $10^{-3}$\\
$E_i,E_0$ & electric field, applied bias field & V\,m$^{-1}$ & $10^{3}$ to $10^{6}$\\
$v_d$ & drift velocity, mobility times bias field & m\,s$^{-1}$ & $1$ to $10^{3}$\\
$j_i$ & particle flux of carriers & m$^{-2}$s$^{-1}$ & $10^{18}$ to $10^{24}$\\
$J_i$ & electric current density, charge times particle flux & A\,m$^{-2}$ & $10^{2}$ to $10^{8}$\\
$S_n$ & Seebeck coefficient, voltage per kelvin, negative for electrons & V\,K$^{-1}$ & $10^{-4}$ to $10^{-3}$\\
$\beta_{\rm th}$ & thermodiffusion coefficient in the carrier equation & m$^{-1}$s$^{-1}$K$^{-1}$ & $10^{16}$ to $10^{21}$\\
$\Pi_{\rm th}$ & Peltier coefficient, temperature times Seebeck & V & $10^{-2}$ to $10^{-1}$\\
$\tau_R$ & recombination time of an excess carrier & s & $10^{-9}$ to $10^{-6}$\\
$E_g$ & band gap energy & J & about $5\times10^{-19}$\\
$\varphi$ & electric potential & V & 0.1 to 10\\
$D_i$ & electric displacement & C\,m$^{-2}$ & $10^{-3}$ to $10^{-1}$\\
$\epsilon_{ik}$ & permittivity of the material & F\,m$^{-1}$ & $10^{-10}$ to $10^{-9}$\\
$e_{kij}$ & piezoelectric constant, charge per area per unit strain & C\,m$^{-2}$ & 0.1 to 2\\
$p_i$ & pyroelectric coefficient & C\,m$^{-2}$K$^{-1}$ & $10^{-6}$ to $10^{-4}$\\
$\alpha_{kl}$ & thermal expansion coefficient & K$^{-1}$ & $10^{-6}$ to $10^{-5}$\\
$\beta_{ij}$ & thermal stress coefficient, stiffness times expansion & Pa\,K$^{-1}$ & $10^{4}$ to $10^{6}$\\
$c_{ijkl}$ & elastic stiffness & Pa & $10^{10}$ to $10^{11}$\\
$\rho$ & mass density & kg\,m$^{-3}$ & $4\times10^{3}$ to $6\times10^{3}$\\
$\rho c_p$ & heat capacity per unit volume & J\,m$^{-3}$K$^{-1}$ & $10^{6}$ to $3\times10^{6}$\\
$k_{\rm th}$ & thermal conductivity & W\,m$^{-1}$K$^{-1}$ & 1 to 100\\
$\kappa$ & thermal diffusivity, conductivity over capacity & m$^2$s$^{-1}$ & $10^{-6}$ to $10^{-5}$\\
$\tau_q$ & relaxation time of the heat flux (MGT) & s & $10^{-12}$ to $10^{-10}$\\
$\tau_T$ & diffusion time across the plate, thickness squared over diffusivity & s & $10^{-13}$ to $10^{-12}$\\
$l_{\rm nl}$ & nonlocal length of the size effect & m & $10^{-9}$ to $10^{-7}$\\
$l_g$ & strain gradient length of the size effect & m & $10^{-9}$ to $10^{-7}$\\
$2h$, $h$ & plate thickness, and half of it & m & $10^{-9}$ to $10^{-3}$\\
$k$ & wavenumber, complex, along the plate & m$^{-1}$ & $10^{4}$ to $10^{8}$\\
$k_r,k_i$ & real part sets wavelength, imaginary part sets decay & m$^{-1}$ & $k_i\ll k_r$\\
$\lambda$ & wavelength, two pi over the real part of the wavenumber & m & $10^{-8}$ to $10^{-4}$\\
$\omega$ & angular frequency, two pi times ordinary frequency & rad\,s$^{-1}$ & $10^{6}$ to $10^{11}$\\
$f$ & ordinary frequency & Hz & $10^{5}$ to $10^{10}$\\
$f(2h)$ & frequency multiplied by thickness, the comparison axis & Hz\,m & 1 to $10^{4}$\\
\bottomrule
\end{tabular}
\end{center}

\section{Glossary}

\textit{Adiabatic} means no heat flows. \textit{Isothermal} means the temperature does not change. They are
the two limits of any thermoelastic problem and they bracket the truth.

\textit{Anisotropy} is having different properties in different directions. A crystal is usually anisotropic;
a glass is not.

\textit{Attenuation} is the weakening of a wave as it travels.

\textit{Boundary condition} is the statement of what happens at an edge. It is not part of the model of the
material.

\textit{Coupling} is any term through which one field affects another. A model with five fields has more than
five couplings, and each must be justified.

\textit{Cutoff} is the lowest frequency at which a mode exists.

\textit{Dispersion} is the dependence of phase velocity on frequency.

\textit{Dissipation} is the irreversible conversion of wave energy into heat.

\textit{Evanescent} describes a disturbance that dies away instead of travelling.

\textit{Mode} is one of the discrete through-thickness patterns a plate can carry.

\textit{Permittivity} is the constant that relates electric displacement to electric field in a material.

\textit{Phase} is the position of a wave in its cycle, measured in radians.

\textit{Polarisation} in a dielectric is the stretching of atoms into dipoles; in a ferroelectric it is the
alignment of those dipoles over a region.

\textit{Pyroelectricity} is polarisation produced by a change in temperature.

\textit{Reciprocity} is the statement that if cause A produces effect B, then cause B produces effect A with
the same constant. It is the reason two coefficients in the plan are really one.

\textit{Relaxation time} is the time a system needs to come back to equilibrium after a small push.

\textit{Strain} is the change in length per unit length. \textit{Stress} is force per unit area.

\textit{Transducer} is a device that converts between electrical and mechanical signals.

\textit{Vector} is a quantity with direction as well as size. \textit{Tensor} is the same idea with two or more
directions attached, which is how stress and stiffness are written.

\section{How to check your own work}

Nobody will look at your numbers as carefully as you will, so the testing has to be yours. Five generic tests
cover most of the ground, and none of them needs a published value to compare against.

\textit{First, units.} Write the unit of every term in every equation before you code it, and make sure the
units agree two by two. This sounds trivial and it catches most sign and algebra errors, because a term that
has the wrong units usually has the wrong sign structure too. Keep the table you write, and put it in the
paper appendix. A referee reads it and relaxes.

\textit{Second, limits.} Every term you add has a limit in which it should disappear. Set the relaxation time
to zero and you must get the classical heat equation. Set the Seebeck coefficient to zero and you must get the
ordinary carrier equation. Set the size lengths to zero and you must get the classical plate. Compute the
reduced case separately, do not just delete a term from your own code, and compare the two.

\textit{Third, symmetry.} Turn the plate upside down. The dispersion curves must be unchanged, for the
symmetric family and the antisymmetric family separately. If a mode that should be symmetric comes out with a
mixture, either the coupling is real or the assembly has a sign error, and telling those apart is worth a day.

\textit{Fourth, switch a term off in two ways.} Set a coupling constant to zero, and separately assemble a
model in which the term was never written. The two answers must agree to six decimal places. This is the test
the plan calls the switching-off test, and it is a test of your code and not of the physics.

\textit{Fifth, compare with a cited source.} When a paper gives a curve, digitise two or three points and
compare. You are looking for the same order and the same trend, not the same digits, because the material
constants will not be identical. If your curve bends the other way, stop and find out why before writing a
sentence about it.

There are three more habits worth having. Keep a log of what you ran and when, because a result you cannot
reproduce is not a result. Plot the thing you are about to quote, because a number read off a table hides a
shape. And when a result is too good, repeat it from a different direction, such as a different formulation of
the same problem, before you believe it.

\end{document}
